\section*{Chapitre \ref{ch:g7:sound-production} --- Le son~: production et propagation}

\begin{solution}{exo:g7:sound-production:1}
La substance qu'un \omterm{def:g3:sound-around-us:sound}{son} traverse --- l'\omterm{def:g2:air-around-us:air}{air}, l'eau, le bois, l'acier,
n'importe quelle foule moléculaire. Le \omterm{def:g7:sound-production:medium}{non-milieu}~: le vide.
\end{solution}

\begin{solution}{exo:g7:sound-production:2}
C'est le relais qui voyage~: un motif filant de resserrements passé
de rang en rang dans la foule. Chaque \omterm{def:g7:states-of-matter:molecule}{molécule} ne fait que gigoter
autour de sa place --- le message se déplace, les messagers restent.
\end{solution}

\begin{solution}{exo:g7:sound-production:3}
Pas de foule, pas de relais~: la cloche vidée laisse les parois
tremblantes de la sonnette bousculer le néant, et aucun message ne se
forme. Les batailles spatiales sont silencieuses pour la même raison
--- les explosions entre les \omterm{def:g4:solar-system:star}{étoiles} n'ont aucune foule pour porter
leur rugissement.
\end{solution}

\begin{solution}{exo:g7:sound-production:4}
L'\omterm{def:g2:air-around-us:air}{air}, environ \qty{340}{m/s}~; l'eau, environ \qty{1500}{m/s}~;
l'acier, environ \qty{5000}{m/s}. Plus le contact entre les \omterm{def:g7:states-of-matter:molecule}{molécules}
est étroit, plus la bousculade se transmet vite~: volantes, en
contact, verrouillées.
\end{solution}

\begin{solution}{exo:g7:sound-production:5}
Les rangs verrouillés de l'acier relaient presque de main en main~:
environ quinze fois l'allure de l'\omterm{def:g2:air-around-us:air}{air}, et le message du rail arrive
donc bien avant le grondement aérien.
\end{solution}

\begin{solution}{exo:g7:sound-production:6}
Le trajet de la \omterm{def:g1:light-and-shadows:source}{lumière} à travers n'importe quel terrain compte pour
instantané à côté de celui du \omterm{def:g3:sound-around-us:sound}{son}~: le choc vu marque le vrai départ
avec une précision bien supérieure à ce que la main sur le
chronomètre peut exploiter.
\end{solution}

\begin{solution}{exo:g7:sound-production:7}
$340 \times 1 = \qty{340}{m}$ pour l'aller-retour~: le mur se tient à
\qty{170}{m}. Répondre \qty{340}{m}, c'est oublier que la seconde a
payé l'aller \emph{et} le retour.
\end{solution}

\begin{solution}{exo:g7:sound-production:8}
Dans l'eau~: $1500 \times 2 = \qty{3000}{m}$ pour l'aller-retour ---
profondeur \qty{1500}{m}. La bévue~: emprunter les \qty{340}{m/s} de
l'\omterm{def:g2:air-around-us:air}{air} pour un voyage sous-marin.
\end{solution}

\begin{solution}{exo:g7:sound-production:9}
$d = 340 \times 2.5 = \qty{850}{m}$. L'éclat du feu d'artifice a
donné un vrai coup de pistolet de départ~; l'avion invisible n'en
offre aucun --- aucun instant d'émission à partir duquel
chronométrer --- et, pire, le rugissement entendu a quitté l'avion il
y a plusieurs secondes~: l'oreille désigne l'endroit où l'avion
\emph{était}.
\end{solution}

\begin{solution}{exo:g7:sound-production:10}
L'oreille immergée est servie par le relais à \qty{1500}{m/s} de
l'eau, l'oreille du bord par celui de l'\omterm{def:g2:air-around-us:air}{air} à \qty{340}{m/s} --- et
l'essentiel du \omterm{def:g3:sound-around-us:sound}{son} du haut-parleur ne franchit jamais la frontière de
la surface pour passer dans l'\omterm{def:g2:air-around-us:air}{air}. \omterm{def:g7:sound-production:medium}{Milieu} plus rapide, message
privé~: la nageuse gagne malgré la distance.
\end{solution}

\begin{solution}{exo:g7:sound-production:11}
Un seul coup, deux relais~: par l'acier et par l'\omterm{def:g2:air-around-us:air}{air}. Acier~:
$3000 \div 5000 = \qty{0.6}{s}$~; \omterm{def:g2:air-around-us:air}{air}~:
$3000 \div 340 \approx \qty{8.8}{s}$ --- le claquement devance la
détonation d'environ \qty{8.2}{s}.
\end{solution}

\begin{solution}{exo:g7:sound-production:12}
Un seul claquement envoie le même instant dans les deux milieux~; à
\qty{750}{m}, l'enregistreur capte l'arrivée par l'eau (sur
l'hydrophone) et l'arrivée par l'\omterm{def:g2:air-around-us:air}{air} (sur le microphone). \omterm{def:g2:measuring-time:duration}{Durées}
attendues~: eau $750 \div 1500 = \qty{0.5}{s}$, \omterm{def:g2:air-around-us:air}{air}
$750 \div 340 \approx \qty{2.2}{s}$ --- un écart d'environ
\qty{1.7}{s}. À l'envers~: connaissant $d$ et la \omterm{def:g7:motion-average-speed:formula}{vitesse} dans l'\omterm{def:g2:air-around-us:air}{air},
l'écart mesuré donne l'instant d'arrivée par l'eau
($t_{\text{air}} - \text{écart}$), et
$v = 750 \div 0.5 = \qty{1500}{m/s}$ --- l'allure de l'eau mesurée
sans avoir jamais chronométré le claquement silencieux lui-même.
\end{solution}

\begin{solution}{pb:g7:sound-production:1}
\textbf{1.} $340 \times 4 = \qty{1360}{m}$ pour l'aller-retour~:
largeur \qty{680}{m}.
\textbf{2.} $340 \times 1.5 = 510$~; la moitié~: \qty{255}{m}.
\textbf{3.} Première paroi~: $340 \times 1 \div 2 = \qty{170}{m}$~;
seconde~: $340 \times 2 \div 2 = \qty{340}{m}$~; le canyon s'étend
sur $170 + 340 = \qty{510}{m}$ le long de cette ligne.
\textbf{4.} La formule est exacte~; la main ne l'est pas~: les
départs et les arrêts humains se dispersent de quelques dixièmes de
seconde --- des dizaines de \omterm{def:g2:measuring-length:units}{mètres} à \qty{340}{m/s}. D'où la vieille
règle~: recommencer, regarder les lectures se serrer, rapporter le
groupe.
\textbf{5.} $1500 \times 0.2 = \qty{300}{m}$ pour l'aller-retour~:
profondeur \qty{150}{m}.
\textbf{6.} $1500 \times 0.4 \div 2 = \qty{300}{m}$.
\textbf{7.} À la frontière air--eau, l'essentiel d'un cri est
renvoyé --- peu de chose entre dans la rivière, et le faible \omterm{ex:g4:how-sound-travels:echo}{écho} du
fond doit retraverser la frontière pour atteindre une oreille dans
l'\omterm{def:g2:air-around-us:air}{air}~: il n'en survit presque rien. Et même en accordant l'\omterm{ex:g4:how-sound-travels:echo}{écho}, des
aller-retours de quelques dixièmes de seconde mettent en échec tout
chronométrage à l'oreille.
\textbf{8.} Le temps d'\omterm{ex:g4:how-sound-travels:echo}{écho} s'allonge brusquement quand le bateau
franchit la cassure --- fond plus profond, aller-retour plus long~:
le sonar dessine la falaise sous-marine.
\textbf{9.} Acier~: $2000 \div 5000 = \qty{0.4}{s}$~; \omterm{def:g2:air-around-us:air}{air}~:
$2000 \div 340 \approx \qty{5.9}{s}$.
\textbf{10.} Les deux relais partent ensemble~: l'écart croît donc
proportionnellement à la distance --- chaque kilomètre ajoute un
retard supplémentaire fixe au relais lent par rapport au rapide.
Mesure l'écart, divise par la différence par kilomètre, et la
distance tombe toute seule --- un seul coup de marteau, sans
chronomètre nécessaire à l'autre bout.
\textbf{11.} Le \omterm{def:g3:sound-around-us:sound}{son} du rail~: le brouillard est une foule de
gouttelettes qui disperse désespérément la \emph{\omterm{def:g1:light-and-shadows:source}{lumière}} mais ne
gêne guère un relais verrouillé dans de l'acier \omterm{def:g3:solids-liquids-gases:solid}{solide} (le \omterm{def:g3:sound-around-us:sound}{son}
aérien survit lui aussi bien au brouillard --- mais le message du
rail est le plus net et le plus rapide des trois).
\textbf{12.} Par exemple~: «~Tout \omterm{ex:g4:how-sound-travels:echo}{écho} paie l'aller-retour~: diviser
par deux avant de croire. Tout \omterm{def:g7:sound-production:medium}{milieu} fixe sa propre allure~: choisir
la \omterm{def:g7:motion-average-speed:formula}{vitesse} avec la foule, jamais par habitude. Et tous les
chronométrages de ce canyon ont commencé par l'œil --- le trajet de
la \omterm{def:g1:light-and-shadows:source}{lumière} compte ici pour instantané, et c'est pourquoi voir tomber
le marteau peut servir de coup de pistolet de départ.~»
\end{solution}
